\section{A3: exact marked-edge divergence and transition-count entropy}
\label{sec:r33-a3}

The stopping count, its occupation measure, and the unscaled edge current are assigned different symbols. A current carries the mark needed by the append map. Slow information is retained as a distributional divergence defect determined by an exact identity; no extra independently controllable current is manufactured by an unproved mesoscopic cancellation.

\subsection{The exact finite-word complex}
Let $H,E$ be standard Borel spaces, $U:H\times E\to H$ a measurable append map, and $R:E\to\mathbb N_{\ge1}$. For a legal word $m_1,\ldots,m_{\ell_N}$ set
\[
 h_k=U(h_{k-1},m_k),\qquad s_k=N^{-1}\sum_{j\le k}R(m_j),\quad s_0=0.
\]
The word may stop when $\sum_{j\le k}R(m_j)\ge N$; then $\ell_N\le N$. Write $S_N=s_{\ell_N}$. Define the unscaled marked edge current, its scaled flow and transition marginal by
\[
 J_N=\sum_{k=1}^{\ell_N}\delta_{(s_k,h_{k-1},m_k)},\quad
 Q_N=J_N/N,\quad \eta_N=(\mathrm{time},\mathrm{tail})_\#Q_N.
\]
The holding occupation is
\[
 L_N(ds,dh)=\sum_{k=1}^{\ell_N}1_{[s_{k-1},s_k)}(s)\,ds\,\delta_{h_{k-1}}(dh).
\]
For a marked measure $J$, define its divergence by
\[
 \langle\operatorname{div}_U J,\varphi\rangle
 =\int[\varphi(s,h)-\varphi(s,U(h,m))]J(ds,dh,dm).
\]
This definition uses exactly the tail $h_{k-1}$ and head $h_k$.

\begin{theorem}[Exact distributional balance]
\label{thm:r33-a3-balance}
For every bounded test $\varphi(s,h)$ continuously differentiable in time,
\[
 \partial_sL_N+\operatorname{div}_U J_N
 =\delta_{(0,h_0)}-\delta_{(S_N,h_{\ell_N})}.
\]
Equivalently the coefficient of $\operatorname{div}_UQ_N$ is $N$, not one. There is no discretization error in this endpoint convention.
\end{theorem}

\begin{proof}
For each $k$ insert $\varphi(s_k,h_{k-1})$ into the endpoint difference. Then
\[
 \varphi(s_k,h_k)-\varphi(s_{k-1},h_{k-1})
 =\int_{s_{k-1}}^{s_k}\partial_s\varphi(s,h_{k-1})ds
    +\varphi(s_k,h_k)-\varphi(s_k,h_{k-1}).
\]
Summing telescopes. Move the two terms to the distributional side, using $\langle\partial_sL,\varphi\rangle=-\int\partial_s\varphi\,dL$. This gives the stated sign and both endpoints. In particular a chronological-only test reduces to the fundamental theorem of calculus.
\end{proof}

Put $D_N=\operatorname{div}_U J_N=N\operatorname{div}_UQ_N$. The identity gives
\[
 |\langle D_N,\varphi\rangle|
 \le2\|\varphi\|_\infty+S_N\|\partial_s\varphi\|_\infty.
\]
Thus on a bounded total-clock set $D_N$ is bounded as a functional on time-$C^1$, history-bounded tests. This is a dual bound, not an asserted $H^{-1}$ compactness theorem on an undefined infinite-dimensional history manifold. If $(L_N,S_N,h_0,h_{\ell_N})$ converges against these tests, $D_N$ converges against them by the identity. In contrast $\operatorname{div}_UQ_N\to0$ only expresses the leading fast stationarity relation.

\subsection{Holding time is not transition count}
The exact renewal relation, in the right-endpoint convention, is
\[
 \int g\,dL_N=N\int\left[\int_0^{R(m)/N}g(s-r,h)dr\right]Q_N(ds,dh,dm).
\]
If $g$ is uniformly Lipschitz in time, its difference from $\int R(m)g(s,h)dQ_N$ is at most
\[
 \frac{\operatorname{Lip}_s(g)}2\sum_k\left(\frac{R(m_k)}N\right)^2.
\]
Consequently the relation $L=\int R(m)Q(\cdot,dm)$ is justified when the maximum scaled holding length tends to zero and the total clock is bounded, together with return-tail uniform integrability for passage of the weighted measure. A macroscopic holding interval must retain the integral over its whole span; it is not equivalent to a point atom weighted by its length.

\subsection{Stopped-chain entropy is an exact chain rule}
Let $P$ be the reference word law with successive kernels $K(h,dm)$, and let $P^q$ use predictable probability kernels $q_k(m)K(h_{k-1},dm)$ with $\int q_kdK=1$. Both laws use the same stopping rule. Assume the relevant entropies are integrable; truncation extends the identity to the value $+\infty$.

\begin{theorem}[Correct stopped likelihood and entropy]
\label{thm:r33-a3-entropy}
On the stopped-word sigma-field,
\[
 \frac{dP^q}{dP}=\prod_{k=1}^{\ell_N}q_k(m_k),\qquad
 H(P^q\mid P)=E^q\sum_{k=1}^{\ell_N}\int q_k\log q_k\,dK.
\]
The normalized cost is divided by $N$ and counted against transitions. It is not integrated against $L_N$.
\end{theorem}

\begin{proof}
Partition the stopped-word space by its length $l\le N$. On that component the joint law factors into its first $l$ conditional kernels, multiplied by the common indicator that the stopping rule first occurs at $l$. Taking the ratio cancels that indicator. The bounded-length likelihood product is normalized, equivalently by successive conditional integration of the stopped martingale. Under $P^q$, the event $k\le\ell_N$ is determined before the $k$th mark. Conditional expectation of its $\log q_k$ term therefore gives the displayed entropy. No optional stopping theorem with an unbounded horizon is needed.
\end{proof}

For a relaxed marked flow $Q(ds,dh,dm)=\eta(ds,dh)q(s,h,m)K(h,dm)$, the corresponding level-2.5 expression is $\int H(qK\mid K)d\eta$. Convexity gives this as a lower bound for mixtures of predictable controls with the same flow. Equality and recovery on a proposed infinite-dimensional state require an additional realization theorem; they do not follow from the chain rule alone.

\subsection{A fully realizable renewal benchmark}
Take iid binary marks with probabilities $(p,1-p)$ and constant return $R=r$. There are exactly $l_N=\lceil N/r\rceil$ selected marks. If their empirical fraction is $a$, the binomial type probability has logarithm
\[
 -l_N\left[a\log\frac ap+(1-a)\log\frac{1-a}{1-p}\right]+O(\log l_N).
\]
At speed $N$ this gives the good rate $r^{-1}H(\mathrm{Bern}(a)\mid\mathrm{Bern}(p))$ on $[0,1]$. Stirling's formula supplies the upper bound; choosing the nearest feasible type supplies recovery for every $a$, including endpoints by continuity. This is a nonempty example in which the transition-clock rate and the exact word construction agree.

\subsection{Conditioning amplitudes are retained}
\begin{proposition}[Subexponential local factors]
\label{prop:r33-a3-conditioning}
Suppose numerator and denominator local coefficients at a regular conditioning datum obey
\[
 A_N(H)=e^{N\Lambda(H)}N^{-d/2}a_N(H),\qquad
 A_N(0)=e^{N\Lambda(0)}N^{-d/2}a_N(0),
\]
where both factors are positive and $\sup_H|\log a_N(H)|=o(N)$. Then
\[
 N^{-1}\log\frac{A_N(H)}{A_N(0)}
 =\Lambda(H)-\Lambda(0)+o(1)
\]
uniformly. At constant order the ratio $a_N(H)/a_N(0)$ remains and need not be one.
\end{proposition}

\begin{proof}
Take logarithms. The common polynomial factor cancels; the two amplitude logarithms are separately $o(N)$. No equality of source-dependent Riesz vectors or terminal matrix amplitudes is asserted.
\end{proof}

\paragraph{Mechanical application still to be established.}
The completed-graph space for singular billiards, complete-past fractional moment, terminal-state Markov-renewal estimates, exponential tightness and recovery for the full chronological state remain application obligations. The exact complex above fixes the indices, sign, mark type and entropy clock, but does not treat every formal divergence-compatible distribution as a realizable word current.
